\section{讨论}
\label{sec:discussion}

\subsection{组合结构为何对演化至关重要}

如表~\ref{tab:strategy} 所示，在 GAIA 上，Global 策略（用于所有主要实验）很早便于 R4 达到 $73.8\%$ 的峰值，随后崩落至 $49.5\%$（峰值--最终差距：$-24.3\%$）。Global 使用了 \ProjectName 的带类型组件，但没有利用这些组件进行隔离；每项编辑都针对所有任务联合评估。在 pass@2 下，即使某项任务的成功概率已经下降，它仍可能被记录为“已解决”，因此低于阈值的回归会绕过跷跷板约束。防止这种崩落需要由可组合性支持的变体隔离：\ProjectName 的组合结构使每项编辑的\emph{预期作用范围}显式化，这是变体隔离把编辑评估限制在目标任务簇、而不是不加区分地在完整任务集上评估的前提（第~\ref{sec:exp-strategy} 节）。

这种关系与类型系统类似：类型不会生成正确程序，却能让错误程序\emph{可检测}。同理，带类型组件不会阻止不良编辑，但能使其作用范围\emph{显式化}，从而支持独立变体。策略比较表明，变体隔离是稳定演化所必需的（缺少变体隔离的 Global 在达到峰值后退化）；没有组合结构，编辑的预期作用范围无从定义，变体隔离也就无法良好定义。不过，组合结构并不保证行为影响有界：$\tau^3$-Bench Telecom 的失败说明，同类型编辑的累积会引发低于阈值的耦合，同时破坏多种对话模式。

\subsection{轨迹丰富度的作用}

\ProjectName 的完整执行轨迹 $\tau$ 提供了标量奖励之外的诊断信息。案例研究（第~\ref{sec:exp-failure} 节）证实了这一点：检测 GAIA 上的奖励黑客（R10 发布，R11 检测到）需要检查改进\emph{如何}发生（利用格式漏洞还是真正检索）；检测 ALFWorld 上的探索不足（R4--R7）则需要追踪编辑类型分布和发布预测准确率。仅凭逐任务二元结果无法恢复这两种信号。

这些观察促成一条设计原则：\textit{反馈信号的丰富度，决定了能够安全开展之演化的复杂程度上限}。如果只有标量奖励，三种病态现象均不可检测：分数变化无法区分奖励黑客与真正改善、探索不足与收敛，也无法区分灾难性遗忘与评估噪声。轨迹结构使每种病态现象都可诊断，前提是存在此前轮次的轨迹用于比较。$\tau^3$-Bench Telecom 的失败展示了这一边界：尽管已有五轮历史轨迹（R2--R6），累积回归仍绕过跷跷板约束，因为没有任何单独编辑越过检测阈值。因此，结构化轨迹记录是检测病态现象的必要条件，却不足以阻止它们：当耦合在逐任务检测阈值以下累积时，轨迹只能在损害发生后记录症状。

\subsection{操作镜像的范围与局限}

强化学习--符号空间镜像是一种设计启发，而不是形式化框架。经典强化学习的收敛保证要求充分探索状态--动作空间；当状态是符号 harness 配置、动作是开放式代码编辑时，这一条件无法满足。在 Global 策略下，GAIA（GPT-5.4）在 15 轮中完全停滞（$\Delta{=}0.0$）；变体隔离消融（第~\ref{sec:exp-strategy} 节）恢复了稳定改善（最终值 $87.4\%=$ 峰值），但没有任何保证表明该结果能延伸到更长时程（变体可能过度专门化），或延伸到任务间依赖阻碍清晰变体分离的任务分布。该镜像也无法预测\textit{哪一种}病态现象会占据主导：$\tau^3$-Bench Telecom 上的灾难性遗忘在 R7 出现；ALFWorld 上的探索不足主导 R4--R7；GAIA 上的奖励黑客直到 R10 才出现。

因此，我们把操作镜像视为设计检查表，而非预测理论：它指出需要防御的失败模式，却不预测其先后顺序、出现时机或相对严重程度。上述三种病态现象具有代表性，但并不穷尽全部情况；其他强化学习现象（例如适配批次偏离部署任务时的分布漂移，以及困难基准上的奖励稀疏）也可能在符号空间表现为类似的失败模式。

\subsection{跨模型家族泛化}

在 ALFWorld 上，Opus 4.6 元智能体为来自三个模型家族的任务智能体演化 harness：
\begin{itemize}
  \item Sonnet 4.6（同一家族）：$83.6\%\to94.8\%$（$+11.2\%$）；
  \item GPT-5.4（不同家族）：$76.9\%\to97.8\%$（$+20.9\%$）；
  \item Qwen3.5-9B（不同家族且更弱）：$53.0\%\to97.0\%$（$+44.0\%$）。
\end{itemize}
反向缩放效应（第~\ref{sec:exp-main} 节）解释了增益大小的次序：增益与基线性能反向相关（Qwen $>$ GPT $>$ Sonnet），而不是取决于与元智能体模型家族的接近程度。三个配置均保持元智能体 Opus 4.6 固定，只改变任务智能体；我们没有评估更弱的元智能体能否实现相当的增益。

一项补充消融（第~\ref{sec:exp-meta-agent} 节）发现，当单智能体 Evolver 与四阶段 AEGIS 流水线使用相同的元智能体模型和基础设施时，二者准确率相当（$86.4\%$ 对 $87.4\%$；$n{=}103$ 时差异处于采样噪声范围内）。这表明，在该元智能体能力水平上，四阶段分解主要提高效率（约少用 $12\%$ token）和可审计性，而不是带来可测量的准确率改善。

\subsection{成本--性能权衡}

如表~\ref{tab:cost-summary} 所示，演化会产生前期计算成本，而该成本可在后续任务调用中摊销。

\begin{table}[!htbp]
\centering
\small
\caption{演化成本汇总。所有主要实验使用 Global（单一 harness）策略；变体隔离一行来自策略消融（第~\ref{sec:exp-strategy} 节）。}
\label{tab:cost-summary}
\resizebox{\textwidth}{!}{%
\begin{tabular}{lccc}
\toprule
\textbf{实验} & \textbf{轮数} & \textbf{总 Token 数} & \textbf{增益} \\
\midrule
GAIA、GPT-5.4（Global） & 15 & 143.7M & $0.0\%$（峰值 $=$ 初始值） \\
GAIA、GPT-5.4（变体隔离，消融） & 15 & 107.8M & $+13.6\%$ \\
ALFWorld、Sonnet 4.6（Global） & 7 & 43.4M & $+11.2\%$ \\
\bottomrule
\end{tabular}}
\end{table}

策略消融（第~\ref{sec:exp-strategy} 节）表明，在 GAIA 上，变体隔离既比 Global 更有效（最终值 $87.4\%$ 对 $49.5\%$），也更高效（107.8M 对 143.7M token）。Token 减少有两个来源：（1）从结构上看，变体隔离下每项编辑只在目标任务簇上评估，而非完整任务集，从而降低逐轮评估成本；（2）Global 下不断退化的基线会使更多候选方案无法通过门控，在从未发布的候选方案上浪费 token。对于演化快速收敛的基准（ALFWorld R4--R7、SWE-bench R2--R3），Global 已经足够，运行时程内没有发生退化。

演化后的 harness 也影响\textit{逐任务推理成本}。在 GAIA 上，逐任务 token 消耗下降约 $25\%$（有针对性的工具选择缩短了轨迹）；在 ALFWorld 上则上升约 $60\%$（任务分解提示词延长了执行）。

部署时，演化后的 harness 是无需元智能体推理的静态产物；演化集之外的任务被路由到在演化集上总体成功率最高的变体。在 GAIA 上，前期 107.8M token 的成本约在 1,300 次调用后摊销完毕（每次调用约节省 83K token）。在 ALFWorld 上，逐任务成本增加；其回报是准确率（$+11.2\%$），而非成本下降。

\subsection{伦理考量}

自演化智能体系统需要显式监督。\ProjectName 提供三种机制：
\begin{enumerate}
  \item \textbf{可审计性}：每项已发布编辑都带有清单和回滚目标；被拒候选方案会连同拒绝原因一起归档。
  \item \textbf{确定性门控}：在 pass@2 下，跷跷板约束会拒绝任何使哪怕一项先前已解决任务退化的编辑。
  \item \textbf{人在回路}：门控层支持对超过可配置风险阈值的编辑进行人工批准（我们的自动化实验未启用该机制）。
\end{enumerate}
$\tau^3$-Bench 的失败（第~\ref{sec:exp-failure} 节）说明了这些机制的局限：连续五次同类型编辑（R2--R6）累积了跷跷板约束未检测到的低于阈值耦合；第六次编辑（R7）触发了可见的 $-14.0\%$ 退化，但此前没有任何单项编辑违反约束。这是逐编辑门控的结构限制：不论先前多少轮在同一约束下表现出表面稳定性，低于阈值的回归都会在未被检测的情况下累积。

\subsection{局限性}

除上文指出的限制外，以下五项约束进一步限定了结果的普适性：
\begin{itemize}
  \item \textbf{没有留出评估。}所有报告的增益都在用于演化的同一任务集上测量。由于我们报告峰值准确率，且在适配集本身上评估，数值同时带有选择偏差与潜在过拟合。对同一分布内未见任务的泛化具有可能性，但尚未经过检验。
  \item \textbf{仅含离散动作空间。}所有实验都使用离散的文本动作空间。我们尚未检验该框架能否扩展到连续动作空间（例如机器人控制）。
  \item \textbf{闭源元智能体。}AEGIS 需要能够进行多文件代码生成、结构化轨迹分析和多步规划的元智能体。接近这一能力水平的开放权重模型（例如 Qwen3.5-72B、Llama-4-Maverick）尚未作为元智能体进行测试。
  \item \textbf{联合控制假设。}协同演化要求同时控制 harness 演化与模型训练。实践中，这两类工作往往分属不同团队或组织；没有跨团队协调时，共享回放缓冲区（第~\ref{sec:coevolution-loop} 节）并不现实。
  \item \textbf{基准覆盖范围。}所有 SWE-bench Verified 实验只使用 55 个任务的子样本，$\tau^3$-Bench 也只评估三个领域（Retail、Airline、Telecom）。相关结论，尤其是反向缩放效应，未必能泛化到任务异构性不同的领域或更大的评估集。
\end{itemize}
